% ch04_espaces_generalises.tex — Espaces M\'etriques G\'en\'eralis\'es
\chapter{Espaces M\'etriques G\'en\'eralis\'es}\label{ch:espaces_generalises}

\section{Introduction}

La notion d'espace m\'etrique, telle que Fr\'echet l'a d\'efinie en 1906, repose sur trois axiomes simples. Mais que se passe-t-il si l'on assouplit l'un d'entre eux~? Si l'in\'egalit\'e triangulaire n'est v\'erifi\'ee qu'\`a une constante multiplicative pr\`es, on obtient un \emph{$b$-espace m\'etrique}. Si la distance d'un point \`a lui-m\^eme n'est pas n\'ecessairement nulle, on entre dans le monde des \emph{espaces m\'etriques partiels}, invent\'es par Steve Matthews en 1994 pour mod\'eliser les domaines de l'informatique. Chacune de ces g\'en\'eralisations na\^it d'un besoin concret --- analyse fractale, s\'emantique des programmes, th\'eorie de la mesure --- et chacune soul\`eve la m\^eme question~: le th\'eor\`eme du point fixe survit-il~?

Ce chapitre montre que la r\'eponse est oui, moyennant des adaptations subtiles, et
pr\'esente les $b$-espaces m\'etriques, les espaces m\'etriques partiels et les
$G$-espaces m\'etriques, avec leurs th\'eor\`emes de points fixes associ\'es.

\section{$b$-Espaces m\'etriques}

\begin{definition}[$b$-espace m\'etrique]
Soit $X$ un ensemble et $s \geq 1$ un r\'eel. Une application
$d : X \times X \to [0, +\infty)$ est une \emph{$b$-m\'etrique} (avec constante $s$)
si pour tous $x, y, z \in X$~:
\begin{enumerate}
    \item $d(x, y) = 0 \Leftrightarrow x = y$;
    \item $d(x, y) = d(y, x)$;
    \item $d(x, z) \leq s \bigl[ d(x, y) + d(y, z) \bigr]$ \quad ($b$-in\'egalit\'e triangulaire).
\end{enumerate}
Le couple $(X, d)$ est alors un \emph{$b$-espace m\'etrique}.
\end{definition}

\begin{remarque}
Pour $s = 1$, on retrouve la d\'efinition classique d'espace m\'etrique. La
constante $s$ mesure le \og{}d\'efaut\fg de l'in\'egalit\'e triangulaire.
\end{remarque}

\begin{exemple}[$b$-m\'etrique naturelle]
Soit $X = \R$ et $d(x, y) = (x - y)^2$. Alors $(X, d)$ est un $b$-espace
m\'etrique avec $s = 2$. En effet, $(x - z)^2 = ((x-y) + (y-z))^2 \leq
2(x-y)^2 + 2(y-z)^2 = 2[d(x,y) + d(y,z)]$ par l'in\'egalit\'e
$(a + b)^2 \leq 2(a^2 + b^2)$.

Cependant, $(X, d)$ n'est pas un espace m\'etrique classique car
$d(0, 2) = 4 > d(0, 1) + d(1, 2) = 1 + 1 = 2$.
\end{exemple}

\begin{exemple}[Espace $\ell^p$ avec $0 < p < 1$]
L'espace $\ell^p$ pour $0 < p < 1$ muni de $d(x, y) = \sum_{n=1}^\infty
\abs{x_n - y_n}^p$ est un $b$-espace m\'etrique avec $s = 2^{1/p - 1}$,
mais ce n'est pas un espace m\'etrique.
\end{exemple}

\begin{definition}[Compl\'etude dans un $b$-espace m\'etrique]
Les notions de suite de Cauchy et de compl\'etude se d\'efinissent comme dans
le cas classique~: $(x_n)$ est de Cauchy si $d(x_m, x_n) \to 0$ quand
$m, n \to \infty$, et $(X, d)$ est complet si toute suite de Cauchy converge.
\end{definition}

\begin{attention}[title=\textbf{Difficult\'es sp\'ecifiques aux $b$-m\'etriques}]
Dans un $b$-espace m\'etrique~:
\begin{itemize}
    \item La topologie induite peut ne pas \^etre m\'etrisable au sens classique;
    \item Une $b$-m\'etrique n'est pas n\'ecessairement continue (la fonction
          $(x, y) \mapsto d(x, y)$ peut ne pas \^etre continue pour la topologie
          induite);
    \item Les boules ouvertes peuvent ne pas \^etre ouvertes pour la topologie
          induite.
\end{itemize}
\end{attention}

\subsection{Principe de contraction dans les $b$-espaces m\'etriques}

\begin{theoreme}[Banach pour les $b$-m\'etriques]\label{thm:banach-b}
Soit $(X, d)$ un $b$-espace m\'etrique complet avec constante $s \geq 1$, et
$T : X \to X$ une application v\'erifiant
\[
    d(Tx, Ty) \leq k \, d(x, y) \quad \forall x, y \in X
\]
avec $0 \leq k < 1/s$. Alors $T$ admet un unique point fixe.
\end{theoreme}

\begin{proof}
Soit $x_0 \in X$ et $x_n = T^n(x_0)$. On a $d(x_{n+1}, x_n) \leq k^n d(x_1, x_0)$.
Pour $m > n$, par la $b$-in\'egalit\'e triangulaire appliqu\'ee r\'ecursivement~:
\begin{align*}
    d(x_n, x_m) &\leq s \, d(x_n, x_{n+1}) + s^2 d(x_{n+1}, x_{n+2}) + \cdots + s^{m-n} d(x_{m-1}, x_m) \\
    &\leq \sum_{i=0}^{m-n-1} s^{i+1} k^{n+i} d(x_1, x_0)
    = s k^n \sum_{i=0}^{m-n-1} (sk)^i \, d(x_1, x_0).
\end{align*}
Puisque $sk < 1$, la s\'erie g\'eom\'etrique converge et
$d(x_n, x_m) \leq \frac{sk^n}{1 - sk} d(x_1, x_0) \to 0$.
Donc $(x_n)$ est de Cauchy et converge vers $x^* \in X$.

En utilisant~: $d(x^*, Tx^*) \leq s[d(x^*, x_{n+1}) + d(Tx_n, Tx^*)]
\leq s[d(x^*, x_{n+1}) + k \, d(x_n, x^*)] \to 0$.

L'unicit\'e suit le m\^eme argument que dans le cas classique.
\end{proof}

\begin{remarque}
La condition $k < 1/s$ est plus restrictive que $k < 1$. Pour $s = 1$ (espace
m\'etrique classique), on retrouve la condition usuelle $k < 1$.
\end{remarque}

\begin{theoreme}[Version am\'elior\'ee, Czerwik 1998]
Si $(X, d)$ est un $b$-espace m\'etrique complet et $T : X \to X$ v\'erifie
$d(Tx, Ty) \leq k \, d(x, y)$ avec $k < 1$ (sans la restriction $k < 1/s$),
alors $T$ a un unique point fixe, \`a condition que $d$ soit continue.
\end{theoreme}

\section{Espaces m\'etriques partiels}

\begin{definition}[Espace m\'etrique partiel]
Un \emph{espace m\'etrique partiel} est un couple $(X, p)$ o\`u
$p : X \times X \to [0, +\infty)$ v\'erifie pour tous $x, y, z \in X$~:
\begin{enumerate}
    \item $x = y \Leftrightarrow p(x, x) = p(x, y) = p(y, y)$;
    \item $p(x, x) \leq p(x, y)$ \quad (petit auto-distance);
    \item $p(x, y) = p(y, x)$ \quad (sym\'etrie);
    \item $p(x, z) \leq p(x, y) + p(y, z) - p(y, y)$ \quad (in\'egalit\'e triangulaire modifi\'ee).
\end{enumerate}
\end{definition}

\begin{remarque}
La diff\'erence essentielle avec une m\'etrique classique est que $p(x, x)$
peut \^etre strictement positif. Ceci est motiv\'e par l'informatique
th\'eorique~: $p(x, x)$ repr\'esente la \og{}taille\fg ou le
\og{}co\^ut\fg de l'information $x$.
\end{remarque}

\begin{exemple}
Sur $X = [0, +\infty)$, $p(x, y) = \max(x, y)$ est une m\'etrique partielle.
On a $p(x, x) = x \geq 0$, avec \'egalit\'e seulement en $x = 0$.
\end{exemple}

\begin{theoreme}[Matthews, 1994]
Soit $(X, p)$ un espace m\'etrique partiel complet et $T : X \to X$ une
contraction, c'est-\`a-dire $p(Tx, Ty) \leq k \, p(x, y)$ avec $k \in [0, 1)$.
Alors $T$ admet un unique point fixe $x^*$, et $p(x^*, x^*) = 0$.
\end{theoreme}

\begin{proof}
La m\'etrique partielle $p$ induit une m\'etrique classique
$d_p(x, y) = 2p(x, y) - p(x, x) - p(y, y)$, et $T$ est aussi une contraction
pour $d_p$. Par le th\'eor\`eme de Banach appliqu\'e \`a $(X, d_p)$ (qui est
complet si $(X, p)$ l'est), $T$ a un unique point fixe $x^*$.
De plus, $p(x^*, x^*) = p(Tx^*, Tx^*) \leq k \, p(x^*, x^*) \leq k^2 p(x^*, x^*)
\leq \ldots$, d'o\`u $p(x^*, x^*) = 0$.
\end{proof}

\section{$G$-Espaces m\'etriques}

\begin{definition}[$G$-espace m\'etrique (Mustafa--Sims)]
Un \emph{$G$-espace m\'etrique} est un couple $(X, G)$ o\`u
$G : X \times X \times X \to [0, +\infty)$ v\'erifie pour tous $x, y, z, a \in X$~:
\begin{enumerate}
    \item $G(x, y, z) = 0 \Leftrightarrow x = y = z$;
    \item $0 < G(x, x, y)$ pour $x \neq y$;
    \item $G(x, x, y) \leq G(x, y, z)$ pour $z \neq y$;
    \item $G(x, y, z) = G(x, z, y) = G(y, z, x) = \ldots$ (sym\'etrie en toutes les variables);
    \item $G(x, y, z) \leq G(x, a, a) + G(a, y, z)$ (in\'egalit\'e rectangulaire).
\end{enumerate}
\end{definition}

\begin{exemple}
Pour tout espace m\'etrique $(X, d)$, on peut d\'efinir
$G(x, y, z) = d(x, y) + d(y, z) + d(x, z)$ ou
$G(x, y, z) = \max\{d(x,y), d(y,z), d(x,z)\}$. Les deux fonctions
d\'efinissent des $G$-m\'etriques.
\end{exemple}

\begin{theoreme}[Mustafa--Sims, 2006]
Soit $(X, G)$ un $G$-espace m\'etrique complet et $T : X \to X$ une application
v\'erifiant
\[
    G(Tx, Ty, Tz) \leq k \, G(x, y, z) \quad \forall x, y, z \in X
\]
avec $k \in [0, 1)$. Alors $T$ admet un unique point fixe.
\end{theoreme}

\begin{proof}
Fixons $x_0$ et posons $x_n = T^n(x_0)$. On a~:
\[
    G(x_n, x_{n+1}, x_{n+1}) = G(Tx_{n-1}, Tx_n, Tx_n) \leq k \, G(x_{n-1}, x_n, x_n)
    \leq k^n \, G(x_0, x_1, x_1).
\]
Pour $m > n$~:
\[
    G(x_n, x_m, x_m) \leq G(x_n, x_{n+1}, x_{n+1}) + G(x_{n+1}, x_m, x_m)
    \leq \sum_{i=n}^{m-1} k^i G(x_0, x_1, x_1) \leq \frac{k^n}{1-k} G(x_0, x_1, x_1).
\]
Donc $(x_n)$ est $G$-Cauchy et converge vers $x^* \in X$.

$G(x^*, Tx^*, Tx^*) \leq G(x^*, x_{n+1}, x_{n+1}) + G(x_{n+1}, Tx^*, Tx^*)
\leq G(x^*, x_{n+1}, x_{n+1}) + k \, G(x_n, x^*, x^*) \to 0$.
Donc $Tx^* = x^*$.
\end{proof}

\begin{remarque}
Jleli et Samet (2012) ont montr\'e que de nombreux r\'esultats de points fixes
dans les $G$-espaces m\'etriques peuvent se d\'eduire de r\'esultats classiques
dans les espaces m\'etriques ordinaires, via la m\'etrique associ\'ee
$d_G(x, y) = G(x, y, y) + G(x, x, y)$. Cette observation a suscit\'e un
d\'ebat sur l'utilit\'e r\'eelle de la notion de $G$-m\'etrique.
\end{remarque}

\section{Espaces m\'etriques rectangulaires}

\begin{definition}[Espace m\'etrique rectangulaire (Branciari)]
Un espace m\'etrique rectangulaire est un couple $(X, d)$ o\`u
$d : X \times X \to [0, +\infty)$ v\'erifie les axiomes (1) et (2) d'un
espace m\'etrique, et o\`u l'in\'egalit\'e triangulaire est remplac\'ee par
l'\emph{in\'egalit\'e rectangulaire}~: pour tous $x, y$ distincts et tous
$u, v \in X \setminus \{x, y\}$ distincts,
\[
    d(x, y) \leq d(x, u) + d(u, v) + d(v, y).
\]
\end{definition}

\begin{theoreme}[Branciari, 2000]
Soit $(X, d)$ un espace m\'etrique rectangulaire complet et $T : X \to X$ une
contraction de Banach ($d(Tx, Ty) \leq k \, d(x,y)$, $k < 1$). Alors $T$ admet
un unique point fixe.
\end{theoreme}

\begin{attention}[title=\textbf{Subtilit\'es}]
La preuve originale de Branciari contenait des lacunes, notamment car dans un
espace rectangulaire, les limites ne sont pas n\'ecessairement uniques et la
topologie peut \^etre pathologique. Des corrections ont \'et\'e apport\'ees
par Sarma et al.\ (2009).
\end{attention}

\section{Espaces m\'etriques modul\'es}

\begin{definition}[Espace m\'etrique modul\'e]
Soit $X$ un ensemble. Une application $\omega : (0, \infty) \times X \times X \to [0, +\infty)$
est un \emph{modulaire m\'etrique} si pour tous $x, y, z \in X$ et tous
$\lambda, \mu > 0$~:
\begin{enumerate}
    \item $\omega(\lambda, x, y) = 0$ pour tout $\lambda > 0$ $\Leftrightarrow$ $x = y$;
    \item $\omega(\lambda, x, y) = \omega(\lambda, y, x)$;
    \item $\omega(\lambda + \mu, x, z) \leq \omega(\lambda, x, y) + \omega(\mu, y, z)$.
\end{enumerate}
\end{definition}

\begin{theoreme}
Soit $X_\omega = \{x \in X : \omega(\lambda, x, x_0) < \infty \text{ pour un } \lambda > 0\}$
l'espace m\'etrique modul\'e associ\'e. Si $X_\omega$ est $\omega$-complet et
$T : X_\omega \to X_\omega$ v\'erifie
$\omega(\lambda, Tx, Ty) \leq k \, \omega(\lambda, x, y)$ pour tout $\lambda > 0$
avec $k < 1$, alors $T$ a un unique point fixe.
\end{theoreme}

\section{Comparaison des g\'en\'eralisations}

\begin{formules}[title=\textbf{Synth\`ese des espaces g\'en\'eralis\'es}]
\begin{center}
\renewcommand{\arraystretch}{1.4}
\begin{tabular}{|l|c|c|}
\hline
\textbf{Espace} & \textbf{Modification} & \textbf{Condition FP} \\
\hline
M\'etrique & -- & $k < 1$ \\
$b$-m\'etrique & $d(x,z) \leq s[d(x,y)+d(y,z)]$ & $k < 1/s$ ou $k < 1$ \\
Partiel & $p(x,x) \geq 0$ & $k < 1$ \\
$G$-m\'etrique & 3 variables & $k < 1$ \\
Rectangulaire & In\'eg.\ rectangulaire & $k < 1$ \\
\hline
\end{tabular}
\end{center}
\end{formules}

\section{Exercices}

\begin{exercice}
Montrer que $d(x,y) = \abs{x - y}^2$ d\'efinit un $b$-espace m\'etrique sur $\R$
avec $s = 2$, mais pas un espace m\'etrique.
\end{exercice}

\begin{exercice}
V\'erifier que $\ell^p$ pour $0 < p < 1$ muni de $d(x,y) = \sum \abs{x_n - y_n}^p$
est un $b$-espace m\'etrique complet et calculer la constante $s$.
\end{exercice}

\begin{exercice}
Soit $(X, p)$ un espace m\'etrique partiel. Montrer que
$d_p(x,y) = 2p(x,y) - p(x,x) - p(y,y)$ est une m\'etrique classique sur $X$.
\end{exercice}

\begin{exercice}
Soit $(X, d)$ un espace m\'etrique. Montrer que $G(x,y,z) = d(x,y) + d(y,z) + d(x,z)$
d\'efinit un $G$-espace m\'etrique. Quel est le lien entre la compl\'etude de
$(X, d)$ et celle de $(X, G)$~?
\end{exercice}

\begin{exercice}
Construire un espace m\'etrique rectangulaire qui ne soit pas un espace m\'etrique.
\textit{Indication~: consid\'erer un ensemble \`a quatre points avec une
m\'etrique bien choisie.}
\end{exercice}

\begin{exercice}
Soit $(X, d)$ un $b$-espace m\'etrique complet avec constante $s = 2$ et
$T : X \to X$ avec $d(Tx, Ty) \leq \frac{1}{3} d(x,y)$. Montrer que $T$ a un
unique point fixe et calculer l'estimation d'erreur a priori.
\end{exercice}
